\chapter{Installation}
\label{ch:installation}

\newthought{Getting \Jarvis\ onto your machine} takes a single \cli{pip} command. This
chapter covers the quick install, the optional extras, how to verify the result, and how to
keep \Jarvis\ in a clean, isolated environment so it never collides with your other Python
work.

\section{What you install}
\label{sec:what-you-install}

\Jarvis\ is distributed on \textsc{PyPI} as a small ecosystem of Python packages
(Table~\ref{tab:packages}). The core package, \code{Jarvis-HEP}, is the only one you must
install explicitly; it pulls in \Operas\ automatically because the built-in quick-starts use
it.

\begin{table}[ht]
  \footnotesize
  \begin{tabular}{@{}lll@{}}
    \toprule
    \textbf{Package} & \textbf{Provides} & \textbf{How it arrives} \\
    \midrule
    \code{Jarvis-HEP}        & core runner, samplers, \textsc{cli} (\cmd{Jarvis}) & you install it \\
    \code{Jarvis-Operas}     & in-process operators (\cmd{jopera})               & pulled in automatically \\
    \code{Jarvis-HEP-Portal} & supporting runtime services                       & pulled in automatically \\
    \code{JarvisPLOT}        & plotting \& flowchart rendering (\cmd{jplot})      & optional, install separately \\
    \bottomrule
  \end{tabular}
  \caption{The \Jarvis\ package ecosystem on \textsc{PyPI}.}
  \label{tab:packages}
\end{table}

\begin{jnote}
Most users never need a \textsc{git} checkout of the framework. Install from \textsc{PyPI}
and work inside a standalone \emph{project} directory (Chapter~\ref{ch:core-concepts}).
\end{jnote}

\section{Requirements}
\label{sec:requirements}

\Jarvis\ runs on \textbf{Python 3.10 and newer}, on Linux and macOS. The ``3.10+'' is a
\emph{supported, continuously-tested range}, not an arbitrary floor: before every release the
test suite runs automatically across the supported 3.10+ interpreters, so any version in that
range is one we have actually validated.\sidenote{If you must use an older Python, \Jarvis\
will not install---upgrade the interpreter or create a dedicated 3.10+ environment as in
Section~\ref{sec:clean-env}.}

\begin{jtip}
Install into a \emph{clean virtual environment}. It keeps \Jarvis's dependency set (NumPy,
SciPy, h5py, PyTorch, and friends) from clashing with other projects, and makes the whole
thing trivial to remove later.
\end{jtip}

\section{Quick install}
\label{sec:quick-install}

For most users, one command is enough:

\begin{shellbox}
python3 -m pip install -U Jarvis-HEP
\end{shellbox}

\noindent The \cli{-U} flag upgrades \Jarvis\ and its dependencies to the latest compatible
versions. This installs the \cmd{Jarvis} executable as a console entry point on your
\code{PATH} (Section~\ref{sec:verify} shows how to confirm where it landed).

\section{Optional extras}
\label{sec:extras}

Beyond the core package, two optional pieces are commonly added: the plotting sibling
package, and any external backends that a particular sampler or calculator needs.

\paragraph{Plotting} lives in the separate \code{JarvisPLOT} package, which provides the
\cmd{jplot} command and is what \cli{Jarvis ... -{}-plot} hands off to
(Chapter~\ref{ch:jarvisplot}):

\begin{shellbox}
python3 -m pip install -U JarvisPLOT
\end{shellbox}

\paragraph{External backends} such as \textsc{MultiNest}, \textsc{Diver}, or CERN
\textsc{root} are independent native programs. Install them following their own
instructions; \Jarvis\ detects and drives them when your task card asks for them. A few
optional Python-based samplers likewise depend on extra packages. In every case the relevant
sampler chapter in Part~\ref{part:samplers} states the exact prerequisite and how to install
it, so you only ever add what your scan actually uses.

\begin{jnote}
A full ``everything'' install is simply the three \textsc{PyPI} packages together:
\begin{shellbox}
python3 -m pip install -U Jarvis-HEP Jarvis-Operas JarvisPLOT
\end{shellbox}
After it you will have the \cmd{Jarvis}, \cmd{jopera}, and \cmd{jplot} commands.
\end{jnote}

\section{A clean environment (recommended)}
\label{sec:clean-env}

Any environment manager works---\code{venv}, \code{conda}, \code{uv}, or \code{pyenv}---as
long as it gives you an isolated 3.10+ interpreter. A \code{pyenv}~+~\code{virtualenv}
recipe:

\begin{shellbox}
pyenv install 3.10.12
pyenv virtualenv 3.10.12 jarvis-hep
pyenv activate jarvis-hep
python3 -m pip install -U Jarvis-HEP
\end{shellbox}

\noindent The plain-\code{venv} equivalent:

\begin{shellbox}
python3 -m venv ~/.venvs/jarvis-hep
source ~/.venvs/jarvis-hep/bin/activate
python3 -m pip install -U Jarvis-HEP
\end{shellbox}

\section{Verifying the install}
\label{sec:verify}

Two commands confirm everything is in place:

\begin{shellbox}
Jarvis --version
Jarvis --help
\end{shellbox}

\noindent \cli{Jarvis -{}-version} prints the \Jarvis\ banner and the runtime version of the
installed package (Figure~\ref{fig:version-banner})---a quick way to check exactly which
release you are running. The coloured dot-matrix on the left is the \Jarvis\ icon; the wordmark
fades from blue to gold, and the final line carries the live version of \emph{your} install.

\begin{figure}[ht]
\centering
\JarvisVersionBanner
\caption[The \cli{Jarvis -{}-version} banner.]{The \cli{Jarvis -{}-version} banner, reproduced
in the project's own palette. The version on the last line is read from your installed package
at runtime; here it shows the release this manual documents.}
\label{fig:version-banner}
\end{figure}

\noindent \cli{Jarvis -{}-help} prints the two entry points (a \textsc{yaml} file, or a
\cli{project} subcommand) and the available options; the full \textsc{cli} is documented in
Chapter~\ref{ch:cli}.

\begin{jwarn}
If your shell reports \cli{Jarvis: command not found}, the executable is not on your
\code{PATH}---almost always because you installed into a different environment than the one
you are using now. Re-activate the environment where you ran \cli{pip install}, and confirm
with \cli{which Jarvis}. Under \code{pyenv}, \cmd{Jarvis} lives in that version's \code{bin/}
directory.
\end{jwarn}

\section{Where to go next}
\label{sec:install-next}

With \Jarvis\ installed and verified, head to Chapter~\ref{ch:quickstart} to create a project
and run your first scan in a couple of minutes.
